\documentclass[11pt]{article}
\usepackage[margin=1in]{geometry}
\usepackage{amsmath,amssymb,amsthm}
\usepackage{booktabs}
\usepackage{enumitem}
\usepackage{hyperref}

\newtheorem{theorem}{Theorem}
\newtheorem{proposition}{Proposition}
\newtheorem{corollary}{Corollary}
\newtheorem{remark}{Remark}
\newtheorem{lemma}{Lemma}

\title{Paper Strategy Memo:\\CKME-Based Conformal Prediction --- Scope, Theory, and Experiments}
\author{Internal Discussion Document}
\date{Updated: 2026-04-10}

\begin{document}
\maketitle

\tableofcontents
\newpage

%% ====================================================================
\section{General Introduction}
\label{sec:intro}
%% ====================================================================

\subsection{What This Memo Is}

This document is an internal strategy memo for developing two papers
from our CKME-CP (Conditional Kernel Mean Embedding + Conformal
Prediction) research project. It serves three purposes:

\begin{enumerate}
  \item \textbf{Organize theoretical directions:} Identify what theorems
    we need, what we can prove, and what gaps remain.
  \item \textbf{Connect theory to experiments:} For each theoretical
    result, specify the experiment that validates it and the current
    status of both.
  \item \textbf{Plan the two-paper split:} Paper~1 (methodology +
    theory) and Paper~2 (adaptive Stage-2 design).
\end{enumerate}

\subsection{The Two-Paper Strategy}

\begin{itemize}
  \item \textbf{Paper~1 (CKME-CP methodology):} A CDF-first conformal
    framework with (a)~a finite-sample conditional coverage bound, and
    (b)~a new ``score homogeneity'' regime where CP absorbs estimation
    bias. Target: JASA T\&M, JRSSB, or Biometrika.
  \item \textbf{Paper~2 (adaptive two-stage design):} Optimal Stage-2
    allocation using $S^0$ scores. Target: Technometrics, IISE
    Transactions, or JCGS.
\end{itemize}

\subsection{What We Have Built}

\begin{enumerate}[label=(\roman*)]
  \item \textbf{CKME-CP method:} A CDF-first conformal framework that
    estimates $\hat F(y \mid x)$ via kernel mean embedding, inverts for
    quantiles, and constructs prediction intervals via split CP with
    score $s(x,y) = |\hat F(y \mid x) - 0.5|$.

  \item \textbf{Adaptive bandwidth and L1/L3 decoupling:} Under
    adaptive $h(x) = c\,\sigma(x)$, pre-CP quantiles (L1) are
    biased and do not converge, yet conditional coverage (L3) converges
    at rate $\sim n^{-2/5}$. This CDF-first-specific decoupling is
    absent from quantile-first methods.

  \item \textbf{Two-stage adaptive design:} Stage~1 trains CKME;
    Stage~2 allocates budget via $S^0$ score (quantile interval width)
    to high-uncertainty regions, then calibrates CP on Stage~2 data.

  \item \textbf{Extensive experiments:}
    Asymptotic consistency (\texttt{exp\_conditional\_coverage}),
    non-Gaussian robustness (\texttt{exp\_nongauss}, 6 simulators),
    one-sided quantile estimation (\texttt{exp\_onesided}),
    score homogeneity mechanism test (\texttt{exp\_score\_homogeneity}),
    and benchmarks against DCP-DR, hetGP, CQR.
\end{enumerate}


%% ====================================================================
\section{Theoretical Landscape}
\label{sec:landscape}
%% ====================================================================

This section reviews the three key existing results on conditional
coverage in conformal prediction, identifies the gap we fill, and
states our central insight.

% ------------------------------------------------------------------
\subsection{Existing Results}
% ------------------------------------------------------------------

\subsubsection{DCP: Qualitative Asymptotic Conditional Coverage}

Chernozhukov, W\"uthrich \& Zhu (PNAS, 2021):

\begin{theorem}[DCP, Theorem~3 (informal)]
If $\hat F(\cdot \mid x) \to F(\cdot \mid x)$ consistently
(in $o_P(1)$ sense, \textbf{no rate required}), then
$P(Y \in \hat C(X) \mid X) = 1 - \alpha + o_P(1).$
Under misspecification, only marginal coverage holds.
\end{theorem}

\textbf{Strength:} Extremely weak assumption (no rate).
\textbf{Limitation:} No finite-sample bound, no rate. ``All-or-nothing.''

\subsubsection{RCP: Quantitative Finite-Sample Bound}

Ben Taieb (ICML, 2025) introduces Rectified Conformal Prediction:
transform score $V(x,y)$ into $\tilde V(x,y) = V(x,y) / \hat\tau(x)$
where $\hat\tau(x) \approx Q_{1-\alpha}(V \mid X = x)$.

\begin{theorem}[RCP, Theorem~4 (informal)]
If $F_{\tilde V \mid X=x} \circ F_{\tilde V}^{-1}$ is $L$-Lipschitz:
\[
  P(Y \in C_\alpha(X) \mid X)
  \ge 1 - \alpha + \varepsilon_\tau(X)
  - \alpha L [F_{\tilde V}(\varphi)]^{n+1}.
\]
\end{theorem}

\textbf{Strength:} Explicit finite-sample bound; score-agnostic.
\textbf{Limitation:} Lipschitz composition assumption; bound may depend
on~$x$; requires extra rectifying transformation.

\subsubsection{CP\textsuperscript{2}: Approximate Conditional Validity}

Plassier, Thiao \& Vayatis (ICLR, 2025) prove:

\begin{theorem}[CP\textsuperscript{2}, Theorem~3.2 (informal)]
\[
  P(Y \in C(X) \mid X = x)
  \ge 1 - \alpha - d_{\mathrm{TV}}(x) - p_{n+1}(x),
\]
where $d_{\mathrm{TV}}(x)$ measures the score distribution mismatch
(total variation between true and estimated score distributions), and
$p_{n+1}(x)$ is a finite-sample calibration error.
\end{theorem}

\textbf{Strength:} Clean decomposition into estimation error and
calibration error.
\textbf{Limitation:} The bound allows $d_{\mathrm{TV}}(x)$ to vary
with~$x$; no analysis of when $d_{\mathrm{TV}}$ is constant.

\paragraph{Critical observation (new, not in CP\textsuperscript{2}).}
When $d_{\mathrm{TV}}(x) = \delta$ (constant for all~$x$), the bound
becomes $P(Y \in C(X) \mid X = x) \ge 1 - \alpha - \delta - p_{n+1}$,
which is \emph{uniform} in~$x$. CP\textsuperscript{2} does not discuss
when or why $d_{\mathrm{TV}}$ would be constant. We show that for
CDF-first scores under adaptive bandwidth, this is exactly what happens
(Section~\ref{sec:score-homog}).

\paragraph{Mapping to our framework.}
In our setting:
\begin{itemize}
  \item $d_{\mathrm{TV}}(x)$: controlled by
    $\varepsilon_n = \sup_{x,y}|\hat F(y \mid x) - F(y \mid x)|$ ---
    our CDF estimation error makes $d_{\mathrm{TV}}(x) \le O(\varepsilon_n)$
    uniformly.
  \item $p_{n+1}(x)$: the split-CP finite-sample calibration error,
    $\approx O(n_{\mathrm{cal}}^{-1/2})$.
\end{itemize}

% ------------------------------------------------------------------
\subsection{The Gap We Fill}
% ------------------------------------------------------------------

\begin{center}
\begin{tabular}{@{}llll@{}}
\toprule
\textbf{Method} & \textbf{Assumption} & \textbf{Conclusion} & \textbf{Rate?} \\
\midrule
DCP   & $\hat F \to F$ (no rate) & $= 1-\alpha + o_P(1)$ & No \\
RCP   & score quantile $\hat\tau(x)$ with rate
      & $\ge 1-\alpha + \varepsilon_\tau(x) - \text{rem.}$ & Yes \\
CP\textsuperscript{2}
      & $d_{\mathrm{TV}}(x)$ bounded
      & $\ge 1-\alpha - d_{\mathrm{TV}}(x) - p_{n+1}(x)$ & Yes \\
\textbf{Ours} & $\varepsilon_n = \sup|\hat F - F|$
      & $\ge 1-\alpha - O(\varepsilon_n) - O(n_{\mathrm{cal}}^{-1/2})$
      & \textbf{Uniform} \\
\bottomrule
\end{tabular}
\end{center}

\textbf{Key distinction from CP\textsuperscript{2}:} Their bound allows
$d_{\mathrm{TV}}(x)$ to vary with~$x$. We show that CDF-first scores
make $d_{\mathrm{TV}}(x)$ \emph{automatically} bounded by $O(\varepsilon_n)$
(uniform in~$x$) via PIT. Furthermore, under adaptive bandwidth, the
score distribution becomes $x$-invariant even when $\varepsilon_n \not\to 0$
--- a regime none of the above results address.

% ------------------------------------------------------------------
\subsection{Our Central Insight: CDF-First Scores Are Naturally Rectified}
\label{sec:central-insight}
% ------------------------------------------------------------------

\begin{quote}
\textbf{If $\hat F(y \mid x) = F(y \mid x)$ (oracle), the abs-median
score $s(x,y) = |F(y \mid x) - 0.5|$ has conditional distribution
$U[0, 0.5]$ for all~$x$. No rectification needed.}
\end{quote}

\paragraph{Why this matters.}
RCP's rectifying transformation aims to make
$Q_{1-\alpha}(\tilde V \mid X = x) = \text{const}$. For CDF-based
scores, PIT gives this for free:
$F(Y \mid X) \mid X = x \sim U[0,1] \Rightarrow |F(Y \mid X) - 0.5|
\mid X = x \sim U[0, \tfrac{1}{2}]$ for all~$x$.

\paragraph{When $\hat F \neq F$ (realistic case).}
$\bigl|\, |\hat F(y \mid x) - 0.5| - |F(y \mid x) - 0.5|\, \bigr|
\le |\hat F(y \mid x) - F(y \mid x)| \le \varepsilon_n.$
The score quantile deviates from $(1-\alpha)/2$ by at most
$O(\varepsilon_n)$, \emph{uniformly in~$x$}.

\paragraph{Two regimes.}
\begin{enumerate}
  \item \textbf{Consistent estimation} ($\varepsilon_n \to 0$):
    Score is asymptotically perfectly rectified $\to$ conditional coverage.
    (Covered by Result~1, Section~\ref{sec:result1}.)

  \item \textbf{Biased but homogeneous} ($\varepsilon_n \not\to 0$ but
    bias is $x$-invariant via $h(x) \propto \sigma(x)$):
    Score distribution is $x$-independent $\to$ CP's single
    $\hat q$ absorbs bias uniformly $\to$ conditional coverage holds.
    (Covered by Result~2, Section~\ref{sec:score-homog}.)
\end{enumerate}

\paragraph{Important subtlety: uniform $d_{\mathrm{TV}}$ alone is
insufficient.}
Making $h \to \infty$ also makes $d_{\mathrm{TV}}(x)$ constant (all
CDFs flatten to 0.5), but the scores lose discriminative power ---
the interval collapses to a global constant-width residual band.
Two conditions are needed simultaneously:
\begin{enumerate}
  \item $d_{\mathrm{TV}}(x) = \text{const}$ (uniform bias), AND
  \item the score retains local $x$-adaptivity (interval width still
    reflects local $\sigma(x)$).
\end{enumerate}
Adaptive $h(x) = c\,\sigma(x)$ with moderate~$c$ achieves both.
The over-smoothed collapse (Proposition~\ref{prop:oversmooth}) violates
condition~2.


%% ====================================================================
\section{Result~1: Finite-Sample Conditional Coverage Bound}
\label{sec:result1}
%% ====================================================================

% ------------------------------------------------------------------
\subsection{Theory}
% ------------------------------------------------------------------

\begin{theorem}[Conditional coverage bound for CDF-based CP (informal)]
\label{thm:cov-bound}
Let $F(y \mid x)$ be continuous in~$y$. Let $\hat F$ be any CDF
estimator with uniform error
$\varepsilon_n = \sup_{x,y}|\hat F(y\mid x) - F(y\mid x)|$.
Let $\hat C_{1-\alpha}(x)$ be the split-CP interval using score
$s(x,y) = |\hat F(y \mid x) - 0.5|$ with $n_{\mathrm{cal}}$
calibration points. Then, for all~$x$:
\[
  \bigl| P\bigl(Y \in \hat C_{1-\alpha}(x) \mid X = x\bigr)
  - (1 - \alpha) \bigr|
  = O(\varepsilon_n) + O(n_{\mathrm{cal}}^{-1/2}),
  \quad \text{uniformly in } x.
\]
\end{theorem}

\begin{remark}[Comparison with existing results]
\begin{itemize}
  \item \emph{vs.\ DCP:} Explicit rate instead of $o_P(1)$; price is
    assuming a rate for $\hat F$.
  \item \emph{vs.\ RCP:} Uniform in~$x$, no Lipschitz assumption, no
    rectification step.
  \item \emph{vs.\ CP\textsuperscript{2}:} We show
    $d_{\mathrm{TV}}(x) \le O(\varepsilon_n)$ automatically for
    CDF-based scores; their result allows $x$-dependent $d_{\mathrm{TV}}$.
  \item \emph{Generality:} Applies to \textbf{any} CDF-based CP
    (CKME, DCP-DR, distributional forests, etc.).
\end{itemize}
\end{remark}

\paragraph{Proof sketch.}
\begin{enumerate}
  \item \textbf{Score perturbation:}
    $|s(x,y) - s^*(x,y)| \le \varepsilon_n$ where
    $s^* = |F(y\mid x) - 0.5|$.
  \item \textbf{Oracle score distribution:}
    $s^*(x, Y) \mid X = x \sim U[0, 0.5]$ for all~$x$ (PIT).
  \item \textbf{Threshold concentration:}
    $|\hat q - (1-\alpha)/2|
    \le \varepsilon_n + O(n_{\mathrm{cal}}^{-1/2})$.
  \item \textbf{Combine:}
    $P(s \le \hat q \mid X = x)
    \ge P(s^* \le \hat q - \varepsilon_n \mid X = x)
    = \min(2(\hat q - \varepsilon_n), 1)$.
    Substituting $\hat q$ yields the result.
\end{enumerate}

\paragraph{Connection to CP\textsuperscript{2} Theorem~3.2.}
Our Result~1 can be viewed as a CDF-first specialization of
CP\textsuperscript{2}'s general framework. The key difference: for
\emph{arbitrary} conformity scores, $d_{\mathrm{TV}}(x)$ can depend
on~$x$ in complex ways. For CDF-based abs-median scores, the PIT
structure guarantees $d_{\mathrm{TV}}(x) \le O(\varepsilon_n)$
uniformly. This is what makes CDF-first methods ``naturally rectified.''

% ------------------------------------------------------------------
\subsection{Experiment: Asymptotic Consistency}
\label{sec:exp-consistency}
% ------------------------------------------------------------------

\textbf{Experiment:} \texttt{exp\_conditional\_coverage} (completed).

This experiment validates Result~1 by measuring L1 (quantile error),
L2 (endpoint error), and L3 (conditional coverage error) as $n$ grows
from 64 to 8192, under both fixed and adaptive bandwidths.

\textbf{Key findings:}
\begin{itemize}
  \item Under fixed bandwidth (CV-tuned): all three layers converge.
    L3 slope $\approx -0.15$ on log-log scale. Validates
    Result~1's prediction that $\varepsilon_n \to 0$ implies
    conditional coverage convergence.
  \item Under adaptive bandwidth ($h = c\,\sigma$): L1 (quantile error)
    does \emph{not} converge, yet L3 (coverage) converges at slope
    $\approx -0.39$, close to the $n^{-2/5}$ minimax rate. This
    validates that Result~1's uniform-$\varepsilon_n$ assumption is
    sufficient but not necessary --- Result~2 (score homogeneity)
    provides the missing explanation.
\end{itemize}

% ------------------------------------------------------------------
\subsection{Status and Next Steps}
% ------------------------------------------------------------------

\begin{itemize}
  \item[$\checkmark$] Proof sketch exists (above).
  \item[$\square$] Formalize into a complete proof with explicit constants.
  \item[$\square$] Determine whether the $O(\varepsilon_n)$ term can be
    tightened when $d_{\mathrm{TV}}(x) = \delta$ (constant). This is
    Direction~A from Section~\ref{sec:directions}.
  \item[$\square$] Connect to Result~3 (CKME consistency rate) for a
    concrete instantiation.
\end{itemize}


%% ====================================================================
\section{Result~2: Score Homogeneity Theorem}
\label{sec:score-homog}
%% ====================================================================

This is the ``biased but homogeneously biased'' regime (Regime~2),
which is \textbf{new to the CP literature}.

% ------------------------------------------------------------------
\subsection{Theory: Core Statement}
% ------------------------------------------------------------------

\begin{theorem}[Score Homogeneity (informal)]
\label{thm:score-homog}
Let $Y \mid x \sim F((y - f(x))/\sigma(x))$ be a location-scale
family. Suppose CKME uses bandwidth $h(x) = c\,\sigma(x)$ for some
constant $c > 0$. Then
\[
  \sup_{x}\, \bigl| P(s(x, Y) \le t \mid x) - G(t) \bigr|
  \to 0 \quad \text{as } n \to \infty,
\]
where $G(t)$ is a limit distribution that does not depend on~$x$.
Consequently, CP's single threshold $\hat q$ provides uniform
conditional coverage.
\end{theorem}

\paragraph{Why this is novel.}
\begin{itemize}
  \item \textbf{New regime:} Existing CP theory has two regimes:
    (a)~consistent $\hat F \to$ conditional coverage, or (b)~no
    consistency $\to$ marginal coverage only. Score homogeneity is a
    \emph{third regime}: estimation is biased, but if the bias
    structure is spatially homogeneous, CP absorbs it.
  \item \textbf{CDF-first specific:} Quantile-first scores
    $\max(\hat q_{\mathrm{lo}} - y, y - \hat q_{\mathrm{hi}})$ do
    not have PIT structure, so bias homogeneity does not translate to
    score homogeneity.
  \item \textbf{Practical implication:} $h$ does not need to vanish;
    it only needs to be \emph{proportional to the noise scale}.
\end{itemize}

% ------------------------------------------------------------------
\subsection{Theory: Mechanism (Effective Bandwidth)}
\label{sec:mechanism}
% ------------------------------------------------------------------

CKME estimates the CDF using a smooth indicator
$\Psi_h(Y_i, y) = \Psi((y - Y_i)/h)$, where $\Psi$ is the logistic
CDF. Under $Y = f(x) + \sigma(x)\varepsilon$, define standardized
$z = (y - f(x))/\sigma(x)$, $Z_i = (Y_i - f(x))/\sigma(x)$:
\[
  \frac{y - Y_i}{h} = \frac{z - Z_i}{h / \sigma(x)}.
\]
The \textbf{effective bandwidth} is
$h_{\mathrm{eff}}(x) = h / \sigma(x)$.

\begin{itemize}
  \item \textbf{Fixed~$h$:} $h_{\mathrm{eff}}(x) = h / \sigma(x)$
    varies with~$x$. Over-smooths where $\sigma(x)$ is small;
    under-smooths where $\sigma(x)$ is large. Score distribution is
    $x$-dependent.
  \item \textbf{Adaptive $h(x) = c\,\sigma(x)$:}
    $h_{\mathrm{eff}}(x) = c$ (constant). In standardized scale:
    $\hat F(y \mid x) \approx \hat G_c(z)$, a common kernel-smoothed
    estimate of the noise CDF. Score
    $|\hat G_c(Z) - 0.5|$ is $x$-independent.
\end{itemize}

\paragraph{Over-smoothed collapse (why $h \to \infty$ fails).}

\begin{proposition}[Over-smoothed collapse]
\label{prop:oversmooth}
In the regime $h \to \infty$ (or fixed~$h$ with
$\sigma(x) \ll h$), the abs-median CP interval collapses to:
\[
  U(x) - L(x) \to 2\,Q_{1-\alpha}(\{|Y_j - \hat f(X_j)|\})
  \quad \text{for all } x,
\]
a global constant (proof: linear Taylor expansion of logistic;
bandwidth cancels in CP inversion).
\end{proposition}

\begin{remark}[Width floor and over-coverage]
\label{rem:width-floor}
The collapsed width
$W^* = 2\,Q_{1-\alpha}(|\sigma(X)\varepsilon|)$ depends on the
\emph{marginal} $\sigma(X)$, not local $\sigma(x)$.
\begin{itemize}
  \item Where $\sigma(x) \ll \mathbb{E}[\sigma(X)]$: interval is
    vastly wider than oracle $\to$ over-coverage (wasted precision).
  \item Where $\sigma(x) \gg \mathbb{E}[\sigma(X)]$: interval is
    narrower than oracle $\to$ coverage deficit.
\end{itemize}
\textbf{Numerical example (exp2):}
$\sigma(x) = 0.01 + 0.2(x-\pi)^2$, $\alpha = 0.1$.
At $x = \pi$: $\sigma(\pi) = 0.01$, oracle width $= 0.033$;
collapsed width $\approx 2.95$ ($90\times$ too wide).
At $x = 0$: oracle width $= 6.53$; collapsed width is $0.45\times$
oracle $\to$ severe under-coverage.

Even with adaptive $h(x) = c\,\sigma(x)$, when $\sigma(x) \to 0$
the CKME CDF has $O(h(x))$ smoothing resolution, creating an
irreducible positive width floor. The global CP correction then adds
further width, causing over-coverage. This is exactly what we observe
in the exp2 width curves at $x \approx \pi$
(see Section~\ref{sec:exp-width}).
\end{remark}

% ------------------------------------------------------------------
\subsection{Theory: Connection to CP\textsuperscript{2} and $d_{\mathrm{TV}}$ Analysis}
\label{sec:dtv-analysis}
% ------------------------------------------------------------------

Our score homogeneity result can be formalized through the lens of
CP\textsuperscript{2}'s Theorem~3.2.

\paragraph{When $d_{\mathrm{TV}}(x) = \delta$ (constant).}
CP\textsuperscript{2}'s bound $P(Y \in C(x) \mid X = x) \ge 1 - \alpha
- d_{\mathrm{TV}}(x) - p_{n+1}(x)$ becomes:
\[
  P(Y \in C(x) \mid X = x) \ge 1 - \alpha - \delta - p_{n+1},
\]
a \emph{uniform} bound independent of~$x$. This is stronger than the
general case: every point $x$ has the same coverage lower bound.

\paragraph{Our contribution: identifying when $d_{\mathrm{TV}}$ is constant.}
CP\textsuperscript{2} does not analyze when $d_{\mathrm{TV}}(x)$ is
$x$-independent. We identify a concrete sufficient condition:

\begin{lemma}[Informal]
Under the location-scale model $Y = f(x) + \sigma(x)\varepsilon$
with $\varepsilon \sim G(\cdot)$ independent of~$x$, if the CDF
estimator uses bandwidth $h(x) = c\,\sigma(x)$, then the CDF-based
score distribution satisfies
$d_{\mathrm{TV}}(F_{s|x}, F_{s|x'}) \to 0$ for all $x, x'$.
Hence $d_{\mathrm{TV}}(x)$ (w.r.t.\ the pooled score distribution) is
asymptotically constant in~$x$.
\end{lemma}

\paragraph{Important subtlety.}
Uniform $d_{\mathrm{TV}}$ alone is \textbf{not sufficient}. Making
$h \to \infty$ also produces constant $d_{\mathrm{TV}}$ (all CDFs
flatten), but the score loses discriminative power
(Proposition~\ref{prop:oversmooth}). Two conditions are needed:
\begin{enumerate}
  \item $d_{\mathrm{TV}}(x)$ is constant (score homogeneity), AND
  \item the score retains local $x$-adaptivity (interval width reflects
    local $\sigma(x)$).
\end{enumerate}
Adaptive $h(x) = c\,\sigma(x)$ with moderate~$c$ achieves both.

% ------------------------------------------------------------------
\subsection{Theory: Location-Scale Generalization (Beyond Gaussian)}
\label{sec:loc-scale}
% ------------------------------------------------------------------

The effective-bandwidth argument (Section~\ref{sec:mechanism}) uses
only the location-scale structure $Y = f(x) + \sigma(x)\varepsilon$,
where $\varepsilon \sim G(\cdot)$ is \emph{any} distribution
independent of~$x$. The key insight:

\begin{quote}
Under $h(x) = c\,\sigma(x)$, the standardized CDF estimate
$\hat G_c(z) = \sum_i c_i(x)\,\Psi((z - Z_i)/c)$ depends on~$x$
only through $c_i(x)$ and the i.i.d.\ $Z_i \sim G(\cdot)$. The
effective bandwidth is $c$ for all~$x$, regardless of whether $G$ is
Gaussian, Student-$t$, Gamma, or any other distribution.
\end{quote}

\textbf{This means:}
\begin{itemize}
  \item Result~2 is not restricted to Gaussian noise.
  \item Any location-scale family (Student-$t$, centered Gamma,
    Gaussian mixture with $x$-independent mixing) satisfies the
    homogeneity condition under adaptive bandwidth.
  \item The proposition can be stated for the full location-scale
    family, strengthening the generality of the result.
\end{itemize}

\paragraph{Provability assessment.}
\begin{itemize}
  \item \textbf{Gaussian case:} Most tractable. $Z \sim N(0,1)$ and
    effective bandwidth~$c$ is constant. CDF bias depends only on~$c$
    and the Gaussian density, not on~$x$. Can be proved rigorously.
  \item \textbf{General location-scale:} Requires technical conditions
    on $G(\cdot)$ (e.g., bounded density, finite moments). The
    argument is essentially the same --- the effective bandwidth
    standardization eliminates $x$-dependence --- but the error bounds
    involve $G$'s regularity.
  \item \textbf{Beyond location-scale:} When $G$ depends on~$x$ (e.g.,
    skewness varies with~$x$), the standardization argument breaks.
    However, our \texttt{exp\_nongauss} experiments show that score
    homogeneity is approximately preserved even for moderate violations
    (Section~\ref{sec:exp-nongauss}).
\end{itemize}

% ------------------------------------------------------------------
\subsection{Experiment: Score Homogeneity Mechanism Test}
\label{sec:exp-homog}
% ------------------------------------------------------------------

\textbf{Experiment:} \texttt{exp\_score\_homogeneity} (completed).

Three-phase design with 3 fixed + 6 adaptive bandwidths:
\begin{center}
\begin{tabular}{@{}llp{3.5cm}l@{}}
\toprule
\textbf{Phase} & \textbf{Simulator} & \textbf{$\sigma$ profile} & \textbf{Purpose} \\
\midrule
1 (mild)    & \texttt{mild}          & ratio $\approx 3\times$   & Theory verification \\
2 (extreme) & \texttt{exp2}          & ratio $\approx 200\times$ & Finite-sample limits \\
3 (non-Gaussian) & \texttt{nongauss\_B2L} & Gamma($k{=}2$)      & Beyond location-scale \\
\bottomrule
\end{tabular}
\end{center}

\textbf{Key findings:}
\begin{enumerate}
  \item \textbf{Adaptive $h$ dramatically improves score homogeneity.}
    Phase~1: $\mathrm{KS}_{\max} = 0.065$ (adaptive $c = 1.5$) vs.\
    0.19 (best fixed). Phase~3 (non-Gaussian): $0.17$ vs.\ $0.44$.
  \item \textbf{Decoupling is empirically strong.}
    Anti-diagonal pattern: fixed~$h$ has low quantile error but high
    coverage gap; adaptive~$h$ has high quantile error but low coverage
    gap. Validates Regime~2.
  \item \textbf{Extreme $\sigma$ ratios reveal finite-sample limits.}
    Phase~2: $\mathrm{KS}_{\max}$ cannot be driven below 0.48 even at
    $c = 5.0$, due to near-zero $\sigma(x)$ regions.
  \item \textbf{$c$-scan reveals Pareto frontier.}
    Optimum typically at $c \approx 1$--$2$ (Gaussian) or
    $c \approx 1.5$--$2$ (non-Gaussian).
\end{enumerate}

% ------------------------------------------------------------------
\subsection{Experiment: Width Curves and Over-Coverage}
\label{sec:exp-width}
% ------------------------------------------------------------------

\textbf{Experiment:} \texttt{exp\_conditional\_coverage}, width curve
analysis (completed).

Width curves (mean $\pm$ 1SD of $U(x) - L(x)$ across macroreps vs.\
$x$, overlaid with oracle width) reveal the mechanism behind
over-coverage at low-noise regions.

\textbf{Key observations:}
\begin{itemize}
  \item \textbf{Fixed $h$ (exp2):} Width stays at $\sim 2$--$3$
    everywhere, including at $x = \pi$ where oracle width is
    $\approx 0.033$. Over-smoothed collapse in action.
  \item \textbf{Adaptive $h$ (exp2):} Width tracks oracle much better,
    but an irreducible floor ($\approx 0.5$--$1.0$) persists at
    $x \approx \pi$ where $\sigma(\pi) = 0.01$. The floor arises
    because the smoothed CDF has $O(h(x))$ resolution; when
    $\sigma(x) \ll h(x)$, the CDF is an S-curve of width $O(h)$
    rather than a step function.
  \item \textbf{Convergence:} Both fixed and adaptive width curves
    converge toward oracle as $n \to \infty$, but adaptive converges
    much faster. The adaptive width floor shrinks as $c$ decreases
    with growing~$n$ (since $h \propto c\,\sigma$ and the optimal~$c$
    can shrink).
\end{itemize}

These width curves provide direct visual evidence for
Proposition~\ref{prop:oversmooth} and the width floor mechanism.

% ------------------------------------------------------------------
\subsection{Practical Direction: Plug-In Adaptive Bandwidth}
\label{sec:plugin-h}
% ------------------------------------------------------------------

Result~2 assumes oracle $\sigma(x)$ for $h(x) = c\,\sigma(x)$. In
practice, $\sigma(x)$ is unknown. Our two-stage design provides a
natural plug-in estimator:

\paragraph{Stage-1 provides replications.}
At each Stage-1 site $x_j$ ($j = 1, \ldots, n_0$), we have $r_0$
replications $Y_{j1}, \ldots, Y_{jr_0}$. The local sample standard
deviation
\[
  \hat\sigma(x_j) = \sqrt{\frac{1}{r_0 - 1}
  \sum_{k=1}^{r_0} (Y_{jk} - \bar Y_j)^2}
\]
is available ``for free'' from Stage-1 data. For arbitrary~$x$,
interpolate via kernel smoothing:
$\hat\sigma(x) = \sum_j w_j(x)\,\hat\sigma(x_j)$ with kernel weights.

\paragraph{Plug-in bandwidth:}
$\hat h(x) = c \cdot \hat\sigma(x)$. This is the \emph{practical
realization} of the adaptive bandwidth in Result~2.

\paragraph{What needs to be proved.}
A ``stability'' result: if
$\sup_x |\hat\sigma(x) - \sigma(x)| / \sigma(x) = o(1)$, then the
score homogeneity property is preserved up to an additional
$O(\|\hat\sigma - \sigma\|_\infty / \sigma_{\min})$ error term. This is
plausible because the effective bandwidth becomes
$h_{\mathrm{eff}}(x) = c\,\hat\sigma(x) / \sigma(x) \approx c + o(1)$.

\paragraph{Why this is feasible.}
\begin{itemize}
  \item The plug-in estimator already exists in our codebase (Stage-1
    replications).
  \item No new experiment infrastructure needed --- just modify the
    bandwidth selection in existing experiments.
  \item The stability argument is a straightforward perturbation of
    Result~2.
  \item Combines theory (stability proposition) with practice (the
    plug-in estimator is practical and interpretable).
\end{itemize}

% ------------------------------------------------------------------
\subsection{Status and Next Steps}
% ------------------------------------------------------------------

\begin{itemize}
  \item[$\checkmark$] Mechanism argument (effective bandwidth) written.
  \item[$\checkmark$] Over-smoothed collapse proved
    (Proposition~\ref{prop:oversmooth}).
  \item[$\checkmark$] Score homogeneity experiment completed (3 phases).
  \item[$\checkmark$] Width curves experiment completed (fixed +
    adaptive).
  \item[$\checkmark$] CP\textsuperscript{2} connection identified.
  \item[$\square$] Formalize Theorem~\ref{thm:score-homog} with
    explicit conditions and proof (Gaussian first, then general
    location-scale).
  \item[$\square$] Prove plug-in stability proposition.
  \item[$\square$] State and prove the Lemma connecting adaptive~$h$
    to uniform $d_{\mathrm{TV}}$.
  \item[$\square$] Implement plug-in $\hat h(x) = c\,\hat\sigma(x)$
    in code and run comparison experiment.
\end{itemize}


%% ====================================================================
\section{Result~3: CKME Consistency Rate}
\label{sec:result3}
%% ====================================================================

Result~1 applies to \emph{any} CDF estimator with uniform error
$\varepsilon_n$. To instantiate it for CKME, we need a convergence
rate for $\sup_{x,y}|\hat F(y\mid x) - F(y\mid x)|$.

% ------------------------------------------------------------------
\subsection{Theory: Existing Rates and Strategy}
% ------------------------------------------------------------------

\paragraph{Known rates.}
\begin{itemize}
  \item RKHS norm:
    $\|\hat\mu_{Y|x} - \mu_{Y|x}\|_{\mathcal{H}}
    = O(\sqrt{\log n / n})$
    [Gr\"unew\"alder et al., ICML 2012].
  \item $L^2$ norm (scalar KRR):
    $O(n^{-2m/(2m+d)})$ under Sobolev smoothness~$m$
    [Caponnetto \& De Vito, 2007].
  \item Sup norm (scalar KRR):
    $O((n^{-1}\log n)^{1/2 - d/(4m)})$ under $m > d/2$
    [recent results, 2024--2025].
  \item Sup norm (vector-valued KRR / CME): \textbf{Not directly
    established} --- this is the gap.
\end{itemize}

\paragraph{Two strategies.}
\begin{enumerate}[label=(\Alph*)]
  \item \textbf{Recommended (safe):} State Result~1 with $\varepsilon_n$
    as abstract input. Cite RKHS-norm rate and note sup-norm follows
    under Sobolev embedding ($\mathcal{H}_Y \hookrightarrow C^0$).
    Defer the gap to future work.
  \item \textbf{Ambitious:} Prove sup-norm rate directly via
    (A)~RKHS-to-sup embedding, or (B)~fix~$y$, treat as scalar KRR,
    take sup over~$y$. Either could be a standalone contribution.
\end{enumerate}

% ------------------------------------------------------------------
\subsection{Status and Next Steps}
% ------------------------------------------------------------------

\begin{itemize}
  \item[$\square$] Decide strategy A vs.\ B after formalizing Results
    1 and~2.
  \item[$\square$] If strategy~B: check whether the smooth indicator
    kernel satisfies the Sobolev embedding condition.
\end{itemize}


%% ====================================================================
\section{Non-Gaussian Robustness}
\label{sec:exp-nongauss}
%% ====================================================================

% ------------------------------------------------------------------
\subsection{Theory: Beyond Location-Scale}
% ------------------------------------------------------------------

The score homogeneity result (Section~\ref{sec:loc-scale}) covers the
full location-scale family $Y = f(x) + \sigma(x)\varepsilon$ with
$\varepsilon \perp x$. Beyond this:

\begin{itemize}
  \item \textbf{Approximate location-scale:} If $Y \mid x$ is
    ``close'' to location-scale (e.g., mild $x$-dependent skewness),
    score homogeneity holds approximately. Can we quantify the
    approximation error?
  \item \textbf{Shape-varying noise:} When the noise shape (not just
    scale) depends on~$x$, the standardization argument breaks. CKME
    can still estimate the CDF, but the score distribution becomes
    $x$-dependent. Result~1 (consistency regime) still applies;
    Result~2 (homogeneity) does not.
\end{itemize}

% ------------------------------------------------------------------
\subsection{Experiment: Six Non-Gaussian Simulators}
% ------------------------------------------------------------------

\textbf{Experiment:} \texttt{exp\_nongauss} (designed, partially run).

Six simulators sharing $f(x) = e^{x/10}\sin x$,
$\sigma_{\mathrm{tar}}(x) = 0.1 + 0.1(x-\pi)^2$, $x \in [0, 2\pi]$:

\begin{center}
\begin{tabular}{@{}llll@{}}
\toprule
& \textbf{Student-$t$} & \textbf{Gamma} & \textbf{Gauss mixture} \\
\midrule
Small & A1S ($\nu = 10$) & B2S ($k = 9$) & C1S ($\pi = 0.02$) \\
Large & A1L ($\nu = 3$)  & B2L ($k = 2$) & C1L ($\pi = 0.10$) \\
\bottomrule
\end{tabular}
\end{center}

\begin{itemize}
  \item A1 (Student-$t$): Location-scale family $\to$ Result~2 applies.
  \item B2 (centered Gamma): Approximately location-scale for large~$k$;
    strong skew for small~$k$ $\to$ tests robustness of homogeneity.
  \item C1 (Gaussian mixture): $x$-independent mixing proportions, but
    bimodal shape $\to$ tests CDF estimation quality.
\end{itemize}

Benchmarks against DCP-DR and hetGP (via R). Config:
$n_0 = 250$, $r_0 = 20$, $n_1 = 500$, $r_1 = 10$.

% ------------------------------------------------------------------
\subsection{Status and Next Steps}
% ------------------------------------------------------------------

\begin{itemize}
  \item[$\checkmark$] Simulators implemented and registered.
  \item[$\checkmark$] CV-tuned hyperparameters pretrained.
  \item[$\square$] Run main experiment ($n_{\mathrm{macro}} \ge 5$).
  \item[$\square$] Analyze: does adaptive $h$ improve coverage
    uniformly across non-Gaussian noise types?
  \item[$\square$] If A1 (Student-$t$) shows strong homogeneity and
    B2L (Gamma) shows weaker: supports our theory's location-scale
    boundary.
\end{itemize}


%% ====================================================================
\section{Strengthening Directions: Brainstorming Summary}
\label{sec:directions}
%% ====================================================================

Based on the theoretical analysis and experimental evidence, we
identified three directions to strengthen the theoretical contribution.

% ------------------------------------------------------------------
\subsection{Direction~A: Tighter Bound When $d_{\mathrm{TV}}$ Is Uniform}
% ------------------------------------------------------------------

\paragraph{Idea.}
When $d_{\mathrm{TV}}(x) = \delta$ for all~$x$,
CP\textsuperscript{2}'s bound can potentially be tightened. The
marginal calibration guarantee of CP already controls the average
coverage; if the deviation from target is \emph{constant} across~$x$,
the two-sided bound may shrink.

\paragraph{Feasibility assessment.}
\begin{itemize}
  \item \textbf{Novelty:} Moderate. The observation that constant
    $d_{\mathrm{TV}}$ yields a uniform bound is somewhat folklore in
    the CP literature, though not explicitly stated in
    CP\textsuperscript{2}.
  \item \textbf{Effort:} Low. A corollary of existing results.
  \item \textbf{Impact:} By itself, insufficient for a paper. Useful
    as a supporting lemma.
\end{itemize}

% ------------------------------------------------------------------
\subsection{Direction~B: Plug-In Adaptive Bandwidth from Stage-1 Replications}
% ------------------------------------------------------------------

\paragraph{Idea.}
Use $\hat\sigma(x)$ from Stage-1 replications to construct
$\hat h(x) = c\,\hat\sigma(x)$. Prove a stability result: if
$\hat\sigma \to \sigma$ uniformly, score homogeneity is preserved
(see Section~\ref{sec:plugin-h}).

\paragraph{Feasibility assessment.}
\begin{itemize}
  \item \textbf{Novelty:} High. No existing work connects Stage-1
    replication-based variance estimation to adaptive bandwidth for
    conformal coverage. This is a practical, actionable contribution.
  \item \textbf{Effort:} Moderate. Perturbation analysis of Result~2
    plus one new experiment (plug-in vs.\ oracle comparison).
  \item \textbf{Impact:} Strong. Makes Result~2 \emph{practical}:
    readers can implement it without knowing $\sigma(x)$.
\end{itemize}

% ------------------------------------------------------------------
\subsection{Direction~C: Non-Gaussian / Location-Scale Generalization}
% ------------------------------------------------------------------

\paragraph{Idea.}
Extend Result~2 from Gaussian location-scale to the full location-scale
family $Y = f(x) + \sigma(x)\varepsilon$, $\varepsilon \sim G(\cdot)$
with $G$ arbitrary. The effective-bandwidth argument already works
(Section~\ref{sec:loc-scale}); the main task is stating and proving
it rigorously.

\paragraph{Feasibility assessment.}
\begin{itemize}
  \item \textbf{Novelty:} Moderate--high. Existing CP theory does not
    address how noise distribution shape interacts with bandwidth
    choice for conditional coverage.
  \item \textbf{Effort:} Low--moderate. The argument is the same as
    the Gaussian case; only the regularity conditions on~$G$ need
    specification.
  \item \textbf{Impact:} Strong when combined with experiments.
    Our 6 non-Gaussian simulators provide empirical validation.
\end{itemize}

% ------------------------------------------------------------------
\subsection{Recommended Path: Combine B + C}
% ------------------------------------------------------------------

\begin{quote}
\textbf{Proposition} (target statement): Under the location-scale model
$Y = f(x) + \sigma(x)\varepsilon$ with $\varepsilon \sim G(\cdot)$
independent of~$x$, if:
\begin{enumerate}
  \item $\hat\sigma(x) \to \sigma(x)$ uniformly (from Stage-1
    replications), and
  \item CKME uses plug-in bandwidth $\hat h(x) = c\,\hat\sigma(x)$,
\end{enumerate}
then the CDF-based conformity score satisfies score homogeneity, and
$P(Y \in C(x) \mid X = x) \ge 1 - \alpha - O(p_{n+1})$
\emph{uniformly in~$x$}, where $p_{n+1}$ is the split-CP calibration
error.
\end{quote}

This combines:
\begin{itemize}
  \item \textbf{Theory:} Direction~A (uniform $d_{\mathrm{TV}}$ bound)
    as a lemma, Direction~B (plug-in stability), Direction~C
    (location-scale generalization) as the main proposition.
  \item \textbf{Experiment:} Score homogeneity mechanism test (validates
    the mechanism), \texttt{exp\_nongauss} (validates non-Gaussian
    robustness), plug-in comparison (validates Direction~B).
  \item \textbf{Practical contribution:} A concrete, implementable
    bandwidth choice rule that provably yields conditional coverage.
\end{itemize}


%% ====================================================================
\section{Paper Organization}
\label{sec:paper-org}
%% ====================================================================

% ------------------------------------------------------------------
\subsection{Paper~1: CKME-CP Methodology}
% ------------------------------------------------------------------

\paragraph{Target:} JASA Theory \& Methods, JRSSB, or Biometrika.

\paragraph{Outline:}
\begin{enumerate}
  \item \textbf{Introduction:} CDF-first vs.\ quantile-first; the
    natural-rectification insight; comparison with DCP/RCP/CP\textsuperscript{2}.
  \item \textbf{Method:} CDF estimation via CKME; abs-median score;
    split-CP construction.
  \item \textbf{Theory:}
    \begin{itemize}
      \item Result~1: Finite-sample conditional coverage bound
        (general, any CDF-based CP).
      \item Result~2: Score Homogeneity Theorem (location-scale +
        adaptive $h$ $\to$ $x$-invariant score; includes plug-in
        stability).
      \item Result~3: CKME consistency rate (instantiates Result~1).
    \end{itemize}
  \item \textbf{Experiments:}
    \begin{itemize}
      \item \texttt{exp\_conditional\_coverage}: L1/L3 decoupling,
        convergence rates, width curves.
      \item \texttt{exp\_score\_homogeneity}: direct mechanism test
        (KS, decoupling, $c$-scan).
      \item \texttt{exp\_nongauss}: non-Gaussian robustness
        (6 simulators, benchmarks).
    \end{itemize}
  \item \textbf{Discussion:} Bandwidth choice guidelines; when
    CDF-first beats quantile-first; limitations.
\end{enumerate}

\paragraph{Narrative flow:}
\begin{center}
\fbox{\parbox{0.90\textwidth}{
\textbf{Observation:} CDF-based scores are naturally rectified via PIT.

$\downarrow$ \emph{What does this buy?}

\textbf{Result~1:} Uniform coverage bound
$O(\varepsilon_n) + O(n_{\mathrm{cal}}^{-1/2})$.

$\downarrow$ \emph{What about the biased regime?}

\textbf{Result~2:} Score homogeneity under adaptive $h \propto \sigma$.
Location-scale family, plug-in bandwidth, conditional coverage without
estimation consistency.

$\downarrow$ \emph{Instantiation}

\textbf{Result~3:} CKME achieves $\varepsilon_n \to 0$, feeding
Result~1. Combined with Result~2, gives a complete picture.

$\downarrow$ \emph{Validation}

\textbf{Experiments:} Mechanism test, convergence rates, width analysis,
non-Gaussian robustness.
}}
\end{center}

\paragraph{Stage-2 design in Paper~1.}
Use simplest design (LHS) for Stage-2 calibration. One sentence:
``More sophisticated adaptive designs are explored in [companion
paper].''

% ------------------------------------------------------------------
\subsection{Paper~2: Adaptive Stage-2 Design}
% ------------------------------------------------------------------

\paragraph{Target:} Technometrics, IISE Transactions, JCGS.

\paragraph{Core message.}
Given a trained CKME model from Stage~1, how to optimally allocate
Stage-2 budget to minimize CP interval width?

\paragraph{What exists:} $S^0$ score, three allocation methods (LHS,
sampling, mixed), empirical comparisons.

\paragraph{What Paper~2 needs:}
\begin{enumerate}
  \item Theory: $S^0$-driven allocation reduces expected width or
    coverage gap (variance reduction / minimax / optimal $\gamma$).
  \item Connection to sequential design literature.
  \item Generality: $S^0$ works with any CDF estimator, not just CKME.
\end{enumerate}

% ------------------------------------------------------------------
\subsection{Why Splitting Does Not Hurt}
% ------------------------------------------------------------------

\begin{center}
\begin{tabular}{@{}lp{5.5cm}p{5.5cm}@{}}
\toprule
\textbf{Concern} & \textbf{Paper~1} & \textbf{Paper~2} \\
\midrule
Cross-dependence &
  Uses simplest Stage-2 design; no dependence on Paper~2. &
  Cites Paper~1 for CKME-CP; standard. \\
\addlinespace
$S^0$ needs CKME &
  Not discussed in detail. &
  $S^0$ via any CDF estimator; CKME is one instantiation. \\
\addlinespace
Reviewer overlap &
  Theory + methodology for CP. &
  Design + allocation for simulation. \\
\bottomrule
\end{tabular}
\end{center}


%% ====================================================================
\section{Relationship Between All Results}
\label{sec:relationships}
%% ====================================================================

\begin{center}
\begin{tabular}{@{}p{2.5cm}p{3.5cm}p{3.5cm}p{3cm}@{}}
\toprule
& \textbf{Assumption} & \textbf{Conclusion} & \textbf{Regime} \\
\midrule
\textbf{Result~1}
  & $\varepsilon_n = \sup|\hat F - F|$ known
  & coverage gap $\le O(\varepsilon_n) + O(n_{\mathrm{cal}}^{-1/2})$
  & Consistent \\
\addlinespace
\textbf{Result~2}
  & $h(x) \propto \sigma(x)$, location-scale
  & score dist.\ $x$-invariant $\to$ cond.\ coverage
  & Biased, homogeneous \\
\addlinespace
\textbf{Result~3}
  & kernel regularity
  & $\varepsilon_n \to 0$ with rate
  & Feeds Result~1 \\
\bottomrule
\end{tabular}
\end{center}

\textbf{Flow:} Result~3 $\to$ Result~1: quantitative conditional
coverage for CKME-CP under consistency. Result~2 covers the
complementary regime where consistency fails but bias is homogeneous.
Together: a complete picture of when and why CKME-CP achieves
conditional coverage.


%% ====================================================================
\section{Timeline and Action Items}
\label{sec:timeline}
%% ====================================================================

\subsection{Phase~1: Theory (Current Priority)}

\begin{enumerate}
  \item Formalize Result~1 (finite-sample bound) with complete proof.
  \item Formalize Result~2 (Score Homogeneity) for Gaussian
    location-scale, then extend to general location-scale (Direction~C).
  \item Prove plug-in stability proposition (Direction~B).
  \item State the uniform-$d_{\mathrm{TV}}$ lemma connecting to
    CP\textsuperscript{2} (Direction~A).
  \item Decide strategy for Result~3 (option~A safe vs.~B ambitious).
\end{enumerate}

\subsection{Phase~2: Experiments}

\begin{enumerate}
  \item Run \texttt{exp\_nongauss} with 6 simulators
    ($n_{\mathrm{macro}} \ge 5$).
  \item Implement plug-in $\hat h(x) = c\,\hat\sigma(x)$ and compare
    against oracle $h(x) = c\,\sigma(x)$ in existing experiments.
  \item (Optional) Group conditional coverage analysis vs.\ RLCP.
\end{enumerate}

\subsection{Phase~3: Writing}

\begin{enumerate}
  \item Write Paper~1 draft.
  \item After Paper~1 submission: start Paper~2.
\end{enumerate}

\end{document}
